\begin{theorem}[Markov’s inequality]
Let \((\Omega,\mathcal F,\mathbb P)\) be a probability space and let
\(X:\Omega\to[0,\infty]\) be a non‑negative random variable.  
For any real number \(a>0\),

\[
\mathbb P\bigl(X\ge a\bigr)\le \frac{\mathbb E[X]}{a}.
\]

Moreover,

\[
\mathbb P\bigl(X\ge a\bigr)=\frac{\mathbb E[X]}{a}
\quad\Longleftrightarrow\quad
\mathbb P\bigl(X\in\{0,a\}\bigr)=1 .
\]
\end{theorem}

\begin{proof}
\begin{enumerate}
\item \textbf{Notation.}  
For the fixed \(a>0\) define the indicator of the event \(\{X\ge a\}\) by
\[
\mathbf 1_{\{X\ge a\}}(\omega)=
\begin{cases}
1, & X(\omega)\ge a,\\[2pt]
0, & X(\omega)< a .
\end{cases}
\]

\item \textbf{Pointwise elementary inequality.}  
For every \(\omega\in\Omega\),
\begin{align}
a\,\mathbf 1_{\{X\ge a\}}(\omega)\le X(\omega). \tag{2.1}
\end{align}
Indeed, if \(X(\omega)\ge a\) the left‑hand side equals \(a\le X(\omega)\);  
if \(X(\omega)<a\) the left‑hand side is \(0\le X(\omega)\) because \(X\ge0\).

\item \textbf{Take expectations.}  
Both sides of \eqref{2.1} are non‑negative, so by monotonicity of expectation
(if \(Y\le Z\) a.s.\ then \(\mathbb E[Y]\le\mathbb E[Z]\)) we obtain
\begin{align}
a\,\mathbb E\!\bigl[\mathbf 1_{\{X\ge a\}}\bigr]\le\mathbb E[X]. \tag{3.1}
\end{align}
Linearity of expectation allows pulling the constant \(a\) out.

\item \textbf{Identify the expectation of an indicator.}  
For any event \(A\), \(\mathbb E[\mathbf 1_A]=\mathbb P(A)\).  Hence
\begin{align}
a\,\mathbb P(X\ge a)\le\mathbb E[X]. \tag{4.1}
\end{align}

\item \textbf{Divide by \(a>0\).}  
Since \(a\) is positive,
\[
\mathbb P(X\ge a)\le\frac{\mathbb E[X]}{a},
\]
which is Markov’s inequality.

\item \textbf{Equality condition.}  
Equality in the final inequality can occur only if equality holds in every
preceding step.

\begin{itemize}
\item \emph{Equality in \eqref{2.1}.}  
Equality a.s.\ in \eqref{2.1} means
\[
a\,\mathbf 1_{\{X\ge a\}}(\omega)=X(\omega)\qquad\text{for }\mathbb P\text{-a.e.\ }\omega,
\]
i.e.
\[
X(\omega)=
\begin{cases}
a, & X(\omega)\ge a,\\[2pt]
0, & X(\omega)< a .
\end{cases}
\tag{6.1}
\]
Thus \(X\) takes only the two values \(0\) and \(a\) (almost surely).

\item \emph{Equality in the monotonicity step \eqref{3.1}.}  
The inequality \(\mathbb E[Y]\le\mathbb E[Z]\) is an equality iff
\(Y=Z\) \(\mathbb P\)-a.s.; hence \eqref{3.1} is an equality precisely when
\eqref{6.1} holds, which is the same condition obtained from \eqref{2.1}.
\end{itemize}

Consequently,
\[
\mathbb P\bigl(X\ge a\bigr)=\frac{\mathbb E[X]}{a}
\quad\Longleftrightarrow\quad
\mathbb P\bigl(X\in\{0,a\}\bigr)=1 .
\]

\item \textbf{Conclusion.}  
For any non‑negative random variable \(X\) and any \(a>0\),

\[
\mathbb P\bigl(X\ge a\bigr)\le\frac{\mathbb E[X]}{a},
\]

with equality exactly when \(X\) assumes only the values \(0\) and \(a\)
almost surely. \(\qed\)
\end{enumerate}
\end{proof}